% !TeX spellcheck = ru_RU
% !TEX root = vkr.tex

\section{Обзор литературы} \label{sec:related}

В настоящем разделе систематизируются основные результаты, относящиеся к предметной области исследования: архитектура больших языковых моделей, методы их параметрически эффективного дообучения, гипотеза поверхностного выравнивания, а также подходы к построению AI-агентов, способных использовать внешние инструменты. Обзор формирует теоретическую базу для последующей постановки эксперимента и интерпретации результатов.

\subsection{Большие языковые модели и архитектура Transformer}

Архитектура Transformer, предложенная в работе Vaswani и др.~\cite{vaswani2017attention}, стала фундаментом современных больших языковых моделей (LLM). Ключевым нововведением этой архитектуры является механизм самовнимания (self-attention), позволяющий модели при обработке каждого токена учитывать контекст всех остальных токенов входной последовательности. В отличие от рекуррентных нейронных сетей, обрабатывающих последовательность последовательно, механизм самовнимания допускает полную параллелизацию вычислений, что существенно ускоряет обучение на больших корпусах текстов. Оригинальная архитектура Transformer включает кодировщик (encoder) и декодировщик (decoder), однако в дальнейшем получили широкое распространение модели, использующие только одну из этих компонент.

Семейство моделей GPT (Generative Pre-trained Transformer), описанное в работе Radford и др.~\cite{radford2019language}, использует архитектуру на основе исключительно декодировщика (decoder-only). Модель обучается авторегрессионно: на каждом шаге она предсказывает следующий токен последовательности, обусловливая предсказание на всех предшествующих токенах. Работа GPT-2~\cite{radford2019language} продемонстрировала, что языковые модели, обученные на больших неразмеченных корпусах, способны решать широкий спектр задач обработки естественного языка без явного обучения с учителем (zero-shot), что обосновало концепцию языкового моделирования как универсального подхода к многозадачному обучению.

Семейство моделей YandexGPT~\cite{yandexgpt2024docs} разработано компанией Яндекс и представляет собой ряд генеративных языковых моделей различного масштаба (YandexGPT Lite, YandexGPT Pro). Эти модели доступны через API Yandex Cloud и поддерживают дообучение (fine-tuning) на пользовательских данных, что делает их пригодной платформой для экспериментального исследования влияния дообучения на специализированные задачи, включая классификацию инструментов в контексте AI-агентов.

\subsection{Методы дообучения больших языковых моделей}

Полное дообучение (full fine-tuning) предобученной языковой модели предполагает обновление всех её параметров на целевом наборе данных. Для моделей, содержащих миллиарды параметров, такой подход сопряжён с рядом практических трудностей: высокие вычислительные затраты, необходимость хранения отдельной копии всех параметров для каждой решаемой задачи, а также риск катастрофического забывания (catastrophic forgetting)~--- явления, при котором модель утрачивает ранее приобретённые знания в ходе обучения на новых данных~\cite{kirkpatrick2017overcoming}.

Проблема катастрофического забывания была систематически исследована в работе Kirkpatrick и др.~\cite{kirkpatrick2017overcoming}, где было показано, что нейронные сети при последовательном обучении на различных задачах склонны к значительной деградации качества на ранее освоенных задачах. Авторы предложили метод Elastic Weight Consolidation (EWC), основанный на регуляризации, которая ограничивает изменение параметров, наиболее важных для ранее решённых задач.

В ответ на указанные ограничения полного дообучения был разработан класс методов, объединённых под общим названием \textbf{Parameter-Efficient Fine-Tuning (PEFT)}. Обзор данного направления представлен в работе Ding и др.~\cite{ding2023parameter}, опубликованной в Nature Machine Intelligence. Авторы систематизируют существующие подходы к параметрически эффективному дообучению и выделяют несколько основных стратегий: добавление небольших обучаемых модулей к фиксированной модели, модификация лишь малого числа существующих параметров, а также обучение дополнительных представлений на входе модели.

Метод \textbf{LoRA} (Low-Rank Adaptation), предложенный в работе Hu и др.~\cite{hu2022lora}, основан на гипотезе о низкоранговой природе обновлений весов при дообучении. Вместо обновления полной матрицы весов $W \in \mathbb{R}^{d \times k}$ метод параметризует приращение через произведение двух матриц низкого ранга:
\begin{equation}
    W' = W + \Delta W = W + B A, \quad B \in \mathbb{R}^{d \times r}, \; A \in \mathbb{R}^{r \times k}, \; r \ll \min(d, k).
    \label{eq:lora}
\end{equation}
При этом исходная матрица $W$ остаётся замороженной, а обучаются только матрицы $A$ и $B$, содержащие в совокупности $r(d+k)$ параметров~--- существенно меньше, чем $dk$ параметров полной матрицы. На практике выбор $r$ в диапазоне от 4 до 64 позволяет обучать менее 1\% от общего числа параметров модели, сохраняя при этом качество, сопоставимое с полным дообучением~\cite{hu2022lora}.

Метод \textbf{Prompt Tuning}, предложенный Lester и др.~\cite{lester2021power}, реализует альтернативную стратегию: вместо модификации параметров модели к входной последовательности токенов добавляются обучаемые непрерывные векторы (soft prompts), которые оптимизируются под конкретную задачу. При этом все параметры самой модели остаются замороженными. Авторы показали, что при увеличении масштаба модели эффективность Prompt Tuning приближается к полному дообучению, при значительно меньшем числе обучаемых параметров.

Подход \textbf{Adapter Layers}, предложенный Houlsby и др.~\cite{houlsby2019parameter}, предполагает встраивание небольших дополнительных нейросетевых модулей (адаптеров) между слоями предобученной модели. Адаптеры обычно имеют архитектуру <<бутылочного горлышка>> (bottleneck): понижение размерности, нелинейная активация, повышение размерности. Обучаются только параметры адаптеров, что позволяет использовать одну замороженную базовую модель для множества задач, переключаясь между ними путём замены адаптеров.

\subsection{Гипотеза поверхностного выравнивания}

Важным вопросом при дообучении языковых моделей является характер изменений, которые оно вносит: приобретает ли модель принципиально новые знания или лишь корректирует форму выдачи уже усвоенной при предобучении информации?

В работе Zhou и др.~\cite{zhou2023lima} была сформулирована \textbf{гипотеза поверхностного выравнивания} (Superficial Alignment Hypothesis). Авторы продемонстрировали, что модель LIMA, дообученная всего на 1000 тщательно отобранных примерах, способна генерировать ответы высокого качества, сопоставимые с результатами моделей, обученных на значительно больших наборах данных с использованием обучения с подкреплением на основе обратной связи от человека (RLHF). Согласно данной гипотезе, практически все фактические знания модель приобретает на этапе предобучения, тогда как дообучение с инструкциями (instruction fine-tuning) преимущественно обучает модель \textit{формату и стилю} взаимодействия с пользователем. Иными словами, выравнивание является <<поверхностным>>: оно затрагивает не глубинные знания модели, а лишь способ их представления.

Эти выводы согласуются с результатами работы Gudibande и др.~\cite{gudibande2023false}, в которой было показано, что модели, дообученные на выходах более мощных проприетарных моделей (так называемая имитация), улучшают стиль генерации, но не приобретают фактических знаний, присущих моделям-учителям. Модели-имитаторы демонстрируют улучшение качества ответов по субъективным оценкам (стиль, связность), однако на объективных бенчмарках, измеряющих фактическую корректность, прирост отсутствует~\cite{gudibande2023false}.

Для настоящего исследования гипотеза поверхностного выравнивания имеет ключевое значение. Если дообучение преимущественно корректирует стиль, а не знания, то выбор инструмента в контексте AI-агента может быть интерпретирован как элемент <<стиля>> генерации~--- устойчивый паттерн оформления ответа, который модель способна усвоить из относительно небольшого числа качественных примеров. Данная интерпретация обосновывает ожидание того, что целенаправленное дообучение на компактном, но качественном наборе данных способно значимо улучшить точность классификации инструментов.

\subsection{AI-агенты и выбор инструментов}

Понятие AI-агента на основе большой языковой модели подразумевает систему, в которой LLM выступает в роли центрального <<мыслящего>> компонента, способного планировать действия, обращаться к внешним инструментам и интерпретировать результаты их выполнения. Обзорная работа Xi и др.~\cite{xi2023rise} систематизирует архитектуры таких агентов и выделяет три ключевых компонента: модуль восприятия (perception), модуль планирования и рассуждений (brain), а также модуль действия (action), включающий взаимодействие с внешними инструментами.

Одним из наиболее влиятельных подходов к организации рассуждений агента является фреймворк \textbf{ReAct}, предложенный Yao и др.~\cite{yao2023react}. ReAct объединяет два процесса: генерацию цепочки рассуждений (reasoning) и выполнение действий (acting) в чередующемся цикле <<мысль~--- действие~--- наблюдение>> (thought--action--observation). Модель последовательно формулирует рассуждение о текущей ситуации, выполняет действие (например, вызов API или поисковый запрос), получает наблюдение (результат действия) и использует его для следующего шага рассуждений. Такая архитектура позволяет агенту адаптивно выбирать инструменты на каждом шаге, исходя из контекста задачи.

В работе Schick и др.~\cite{schick2023toolformer} был предложен подход \textbf{Toolformer}, в котором языковая модель самостоятельно обучается вставлять вызовы внешних инструментов (калькулятор, поисковая система, переводчик и др.) в генерируемый текст. Обучение осуществляется в формате self-supervised: модель генерирует кандидатные вызовы инструментов, фильтрует их по критерию снижения перплексии и дообучается на отфильтрованных данных. Toolformer продемонстрировал, что модели способны научиться определять, \textit{когда} и \textit{какой} инструмент следует использовать, без явной разметки.

Работа Patil и др.~\cite{patil2023gorilla} непосредственно связана с задачей настоящего исследования. Авторы предложили модель \textbf{Gorilla}, дообученную для корректного вызова API из обширного каталога. Gorilla обучалась на наборе данных, сгенерированном с помощью Self-Instruct, и продемонстрировала способность точно генерировать вызовы API, превосходя по этому показателю модели GPT-4 без специализированного дообучения. Данный результат подтверждает гипотезу о том, что целенаправленное дообучение на задачу выбора и вызова инструментов способно существенно повысить качество классификации.

Обзорная работа Qin и др.~\cite{qin2024toollearning} систематизирует направление обучения моделей использованию инструментов (tool learning). Авторы выделяют четыре ключевых этапа взаимодействия модели с инструментами: понимание задачи, планирование использования инструментов, непосредственный вызов инструмента и обработка его результата. Особый интерес представляет выделенный авторами этап \textit{выбора инструмента} (tool selection)~--- задача классификации, в которой модель должна из доступного набора инструментов выбрать наиболее подходящий для решения текущей подзадачи. Именно эта задача является предметом настоящего исследования.

\subsection{Синтетическая генерация обучающих данных}

Подготовка качественных обучающих данных для дообучения языковых моделей является ресурсоёмкой задачей. Метод \textbf{Self-Instruct}, предложенный Wang и др.~\cite{wang2023selfinstruct}, предлагает подход к автоматической генерации обучающих примеров с использованием самой языковой модели. Процедура включает итеративную генерацию инструкций, входных данных и эталонных ответов с последующей фильтрацией некачественных и дублирующихся примеров. Авторы продемонстрировали, что дообучение модели на синтетически сгенерированных данных приводит к значимому улучшению качества следования инструкциям, приближаясь к результатам моделей, обученных на данных, размеченных вручную.

Результаты Self-Instruct~\cite{wang2023selfinstruct} согласуются с выводами LIMA~\cite{zhou2023lima}: если выравнивание является поверхностным и затрагивает преимущественно формат генерации, то для его достижения достаточно относительно небольшого числа качественных примеров, даже если эти примеры получены синтетическим путём. Для задачи классификации инструментов это означает, что эффективный обучающий набор может быть сформирован с помощью LLM-генерации при условии обеспечения разнообразия сценариев и корректности разметки. Модель Gorilla~\cite{patil2023gorilla} является практическим подтверждением данного подхода: её обучающий набор был получен методом Self-Instruct.

\subsection{Выводы обзора}

Проведённый обзор литературы позволяет сформулировать следующие ключевые наблюдения.

Во-первых, архитектура Transformer~\cite{vaswani2017attention} и авторегрессионные decoder-only модели~\cite{radford2019language} являются технологическим фундаментом современных LLM, включая семейство YandexGPT~\cite{yandexgpt2024docs}, используемое в настоящей работе.

Во-вторых, методы параметрически эффективного дообучения, в особенности LoRA~\cite{hu2022lora}, позволяют адаптировать модели к специализированным задачам при минимальных вычислительных затратах и сниженном риске катастрофического забывания~\cite{kirkpatrick2017overcoming, ding2023parameter}.

В-третьих, гипотеза поверхностного выравнивания~\cite{zhou2023lima, gudibande2023false} обосновывает ожидание того, что дообучение способно значимо улучшить паттерны генерации (включая выбор инструмента), не требуя при этом масштабных обучающих наборов.

В-четвёртых, задача выбора инструмента в контексте AI-агентов~\cite{xi2023rise, yao2023react, qin2024toollearning} является активной областью исследований, а работы Toolformer~\cite{schick2023toolformer} и Gorilla~\cite{patil2023gorilla} демонстрируют эффективность дообучения для её решения.

Вместе с тем, \textbf{открытым остаётся вопрос} о том, каким образом дообучение модели масштаба YandexGPT Lite влияет на качество классификации инструментов в реальной агентной системе, каков минимальный объём обучающих данных, достаточный для значимого улучшения, и в какой мере наблюдаемый эффект согласуется с гипотезой поверхностного выравнивания. Настоящая работа направлена на экспериментальное исследование именно этих вопросов.
